Para el desarrollo del proyecto se seguirá un enfoque de \textbf{desarrollo iterativo e incremental} (\textit{Iterative and Incremental Development}, \textit{IID}), adaptado al contexto de un \textbf{Trabajo Fin de Grado} desarrollado por un único estudiante. Este tipo de proceso se basa en \textbf{construir el sistema de forma progresiva}, incorporando incrementos funcionales y revisando las decisiones adoptadas a partir de los resultados obtenidos en cada iteración \citep{larman2003iid}. En este proyecto, dicho enfoque permitirá dividir el trabajo en fases progresivas, revisar los resultados obtenidos en cada etapa y ajustar las decisiones técnicas y funcionales conforme avance el desarrollo.

Aunque no se aplicará una \textit{metodología ágil formal} en sentido estricto, el proceso incorporará prácticas habituales de los enfoques iterativos, como la \textbf{revisión continua del alcance}, la \textbf{validación progresiva de requisitos}, la \textbf{priorización de funcionalidades} y la \textbf{generación incremental de entregables}. Esta adaptación resulta adecuada para un proyecto académico con un cliente o usuario experto que aporta información sobre el dominio del problema, pero en el que la ejecución técnica recaerá principalmente sobre el autor del trabajo.

El ciclo de vida previsto se divide en las siguientes fases:

\begin{itemize}
    \item \textbf{Inicio y definición del proyecto:} se identificará el problema, se delimitará el alcance inicial y se establecerán los objetivos generales del sistema.
    \item \textbf{Captura y análisis de requisitos:} se realizarán reuniones y se analizará la situación actual del distribuidor, identificando las necesidades principales del negocio y los procesos susceptibles de digitalización.
    \item \textbf{Estudio de alternativas y selección tecnológica:} se revisarán las herramientas conocidas por el distribuidor y las tecnologías disponibles para el desarrollo de la solución.
    \item \textbf{Análisis y modelado del sistema:} se definirán los módulos funcionales, requisitos funcionales, requisitos no funcionales, requisitos de información, reglas de negocio y modelos conceptuales.
    \item \textbf{Diseño del sistema:} se definirá la arquitectura, el diseño de datos, la estructura de la interfaz y la organización técnica de la aplicación.
    \item \textbf{Construcción incremental:} se implementará la base funcional del sistema, priorizando la autenticación, la gestión de usuarios, el control de roles y la estructura necesaria para futuras ampliaciones.
    \item \textbf{Pruebas y validación:} se comprobará el funcionamiento de los componentes construidos, se validarán manualmente los flujos principales y se revisará la adecuación del sistema respecto a los requisitos definidos.
    \item \textbf{Documentación y cierre:} se elaborarán la memoria, los manuales, los anexos y la documentación técnica necesaria para reproducir y continuar el desarrollo del sistema.
\end{itemize}

La elección del enfoque \textit{IID} se justifica por la amplitud del sistema planteado. El proyecto aborda un dominio con múltiples áreas funcionales, como la \textbf{gestión comercial}, clientes, catálogo, pedidos, entregas, cobros, formaciones e integraciones externas. Por ello, se priorizará la construcción de un \textbf{núcleo técnico modular y ampliable}, sobre el que puedan incorporarse de forma progresiva el resto de módulos definidos en el análisis.

Durante el desarrollo, las decisiones técnicas y funcionales se revisarán de forma continua atendiendo a tres criterios principales: \textbf{adecuación a las necesidades reales del distribuidor}, \textbf{viabilidad técnica} dentro del tiempo disponible y \textbf{coherencia con los objetivos académicos} del Trabajo Fin de Grado. Esta forma de trabajo permitirá mantener el proyecto dentro de un alcance asumible, sin perder de vista la evolución futura del sistema.
